\documentclass[10pt]{extarticle}
\usepackage[a4paper,margin=1.5cm]{geometry}
\usepackage[T1]{fontenc}
\usepackage{amsmath,amssymb}
\usepackage{graphicx}
\usepackage[export]{adjustbox}
\usepackage{multicol}
\usepackage{enumitem}
\usepackage{booktabs,array,tabularx}
\setlength{\columnsep}{0.7cm}
\setlength{\parindent}{0pt}
\setlength{\parskip}{3pt}
\pagestyle{empty}
\makeatletter
\newcommand{\qfigdef}[4]{\@namedef{qf@file@#1}{#2}\@namedef{qf@w@#1}{#3}\@namedef{qf@wide@#1}{#4}}
\newcommand{\qf@check}[1]{\@ifundefined{qf@file@#1}{\PackageError{score2}{Unknown figure id #1}{}}{}}
\newcommand{\qf@img}[4][]{\includegraphics[width=#2\dimexpr\@nameuse{qf@w@#4}pt\relax,max width=#3,max height=0.45\textheight,#1]{\@nameuse{qf@file@#4}}}
\newif\ifqf@full
\newif\ifqf@boxed
\newsavebox{\qf@box}
\newcommand{\qf@open}{\global\qf@fulltrue\par\global\qf@boxedtrue
  \setbox\qf@box\vbox\bgroup\hsize\textwidth\linewidth\textwidth}
\newcommand{\qf@emit}[1]{\par\egroup\global\qf@boxedfalse
  \ifdim\dimexpr\ht\qf@box+\dp\qf@box\relax<#1\textheight\noindent\usebox{\qf@box}\par\else\unvbox\qf@box\fi}
\newcommand{\qfullsplit}{\ifqf@boxed\qf@emit{0.85}\fi}
\newcommand{\qf@close}{\ifqf@boxed\qf@emit{0.45}\fi\global\qf@fullfalse}
\newcommand{\qfullbegin}{\end{multicols}\qf@open}
\newcommand{\qfullend}{\qf@close\begin{multicols}{2}}
\newcommand{\qfullstart}{\qf@open}
\newcommand{\qfulldone}{\qf@close}
\newcommand{\qfig}[2][1]{\qf@check{#2}%
  \ifqf@full\qf@blk{#1}{#2}\else
  \ifnum\@nameuse{qf@wide@#2}=1 \expandafter\@firstoftwo\else\expandafter\@secondoftwo\fi
  {\end{multicols}\begin{center}\qf@img{#1}{\textwidth}{#2}\end{center}\begin{multicols}{2}}%
  {\qf@blk{#1}{#2}}\fi}
\newcommand{\qf@blk}[2]{\par\smallskip{\centering\qf@img{#1}{\linewidth}{#2}\par}\smallskip}
\newcommand{\qfigi}[2][1]{\qf@check{#2}\qf@img[valign=c]{#1}{\linewidth}{#2}}
\makeatother
\newcommand{\papersec}[1]{\par\bigskip{\centering\Large\bfseries PAPER-#1\par}\medskip}
\newcommand{\partsec}[2]{\par\medskip\noindent{\large\bfseries PART-#1: #2}\par\smallskip}
\newcommand{\sect}[1]{\par\smallskip\noindent\textbf{#1}\par}
\newcommand{\qn}[1]{\par\smallskip\noindent\textbf{#1.}~}
\newcommand{\paraq}[2]{\par\smallskip\noindent\textbf{Paragraph for Questions #1 and #2}\par}
\newenvironment{qoptions}{\begin{itemize}[leftmargin=1.8em,itemsep=1pt,topsep=2pt,label={}]}{\end{itemize}}
\newcommand{\opt}[2]{\item[(#1)] #2}
% figure registry (generated by score2_tools.py finalize)
\qfigdef{p1_phy_s1_q01_1}{SRG (Enthusiast) - Test 38 (15 Sep 2024)-img/p1_phy_s1_q01_1.png}{92.3}{0}
\qfigdef{p1_phy_s1_q02_1}{SRG (Enthusiast) - Test 38 (15 Sep 2024)-img/p1_phy_s1_q02_1.png}{85.5}{0}
\qfigdef{p1_phy_s1_q03_1}{SRG (Enthusiast) - Test 38 (15 Sep 2024)-img/p1_phy_s1_q03_1.png}{103.0}{0}
\qfigdef{p1_phy_s1_q04_1}{SRG (Enthusiast) - Test 38 (15 Sep 2024)-img/p1_phy_s1_q04_1.png}{124.6}{0}
\qfigdef{p1_phy_s1_q05_1}{SRG (Enthusiast) - Test 38 (15 Sep 2024)-img/p1_phy_s1_q05_1.png}{100.9}{0}
\qfigdef{p1_phy_s1_q06_1}{SRG (Enthusiast) - Test 38 (15 Sep 2024)-img/p1_phy_s1_q06_1.png}{87.5}{0}
\qfigdef{p1_phy_s1_q07_1}{SRG (Enthusiast) - Test 38 (15 Sep 2024)-img/p1_phy_s1_q07_1.png}{100.1}{0}
\qfigdef{p1_phy_s1_q08_1}{SRG (Enthusiast) - Test 38 (15 Sep 2024)-img/p1_phy_s1_q08_1.png}{85.9}{0}
\qfigdef{p1_phy_s1_q09_1}{SRG (Enthusiast) - Test 38 (15 Sep 2024)-img/p1_phy_s1_q09_1.png}{67.1}{0}
\qfigdef{p1_phy_s2_q01_1}{SRG (Enthusiast) - Test 38 (15 Sep 2024)-img/p1_phy_s2_q01_1.png}{82.1}{0}
\qfigdef{p1_phy_s2_q02_1}{SRG (Enthusiast) - Test 38 (15 Sep 2024)-img/p1_phy_s2_q02_1.png}{97.3}{0}
\qfigdef{p1_phy_s2_q05_1}{SRG (Enthusiast) - Test 38 (15 Sep 2024)-img/p1_phy_s2_q05_1.png}{47.5}{0}
\qfigdef{p1_phy_s2_q06_2}{SRG (Enthusiast) - Test 38 (15 Sep 2024)-img/p1_phy_s2_q06_2.png}{176.6}{0}
\qfigdef{p1_phy_s2_q07_1}{SRG (Enthusiast) - Test 38 (15 Sep 2024)-img/p1_phy_s2_q07_1.png}{64.6}{0}
\qfigdef{p1_phy_s2_q08_1}{SRG (Enthusiast) - Test 38 (15 Sep 2024)-img/p1_phy_s2_q08_1.png}{100.6}{0}
\qfigdef{p1_phy_s2_q09_1}{SRG (Enthusiast) - Test 38 (15 Sep 2024)-img/p1_phy_s2_q09_1.png}{155.8}{0}
\qfigdef{p1_phy_s2_q10_1}{SRG (Enthusiast) - Test 38 (15 Sep 2024)-img/p1_phy_s2_q10_1.png}{111.7}{0}
\qfigdef{p1_chem_s1_q01_1}{SRG (Enthusiast) - Test 38 (15 Sep 2024)-img/p1_chem_s1_q01_1.png}{63.7}{0}
\qfigdef{p1_chem_s1_q01_2}{SRG (Enthusiast) - Test 38 (15 Sep 2024)-img/p1_chem_s1_q01_2.png}{119.6}{0}
\qfigdef{p1_chem_s1_q02_A}{SRG (Enthusiast) - Test 38 (15 Sep 2024)-img/p1_chem_s1_q02_A.png}{55.6}{0}
\qfigdef{p1_chem_s1_q09_1}{SRG (Enthusiast) - Test 38 (15 Sep 2024)-img/p1_chem_s1_q09_1.png}{282.7}{1}
\qfigdef{p1_chem_s2_q03_1}{SRG (Enthusiast) - Test 38 (15 Sep 2024)-img/p1_chem_s2_q03_1.png}{134.4}{0}
\qfigdef{p1_chem_s2_q03_2}{SRG (Enthusiast) - Test 38 (15 Sep 2024)-img/p1_chem_s2_q03_2.png}{300.1}{1}
\qfigdef{p1_chem_s2_q04_1}{SRG (Enthusiast) - Test 38 (15 Sep 2024)-img/p1_chem_s2_q04_1.png}{209.2}{0}
\qfigdef{p1_chem_s2_q08_1}{SRG (Enthusiast) - Test 38 (15 Sep 2024)-img/p1_chem_s2_q08_1.png}{111.7}{0}

\begin{document}
\begin{center}{\Large\bfseries SRG (Enthusiast) --- Test 38 (15 Sep 2024)}\end{center}
\begin{multicols}{2}
\partsec{1}{PHYSICS}
\sect{Section-I}
\qn{1} A circular wire loop (radius $r$, resistance $R$) encloses a region of uniform magnetic field, $B$, perpendicular to its plane. The field increases linearly with time ($B=\alpha t$). An ideal voltmeter (infinite internal resistance) is connected between points P and Q.
\qfig{p1_phy_s1_q01_1}
\begin{qoptions}
\opt{A}{The current in the loop is $\dfrac{\pi\alpha r^2}{2R}$}
\opt{B}{The reading of voltmeter is $\dfrac{\alpha r^2}{2}$}
\opt{C}{The reading of voltmeter is $\alpha r^2$}
\opt{D}{The current in loop is $\dfrac{\pi\alpha r^2}{R}$}
\end{qoptions}
\qn{2} A conducting rod of length $\ell$ slides in a vertical plane against two insulated perpendicular surfaces as shown. Find the induced emf across the rod if its end A is pulled with $\vec v=v_0\hat i$. Outward uniform magnetic field of induction $B$ is applied.
\qfig{p1_phy_s1_q02_1}
\begin{qoptions}
\opt{A}{$\varepsilon=\dfrac{B\ell V_0}{2}$ when $\theta=30^\circ$}
\opt{B}{$\varepsilon=B\ell V_0$ when $\theta=30^\circ$}
\opt{C}{$\varepsilon=0$ when $\theta=45^\circ$}
\opt{D}{$\varepsilon=\dfrac{B\ell V_0}{\sqrt2}$ when $\theta=45^\circ$}
\end{qoptions}
\qn{3} Given figure is a device to fire (project) a body (sliding bar) with a great speed in the order of km/s by an Ampere's force on the conducting sliding bar offered by an impulsive current sent by the charged capacitor after closing the key in the inward magnetic field.\par
If $Q_0$ is the initial charge of the capacitor $C$, $M$ = mass of the sliding bar and $B$ = magnetic field (produced by the closely spaced parallel rails assumed to be uniform), answer the following questions: (at $t=0$, K is closed)
\qfig{p1_phy_s1_q03_1}
\begin{qoptions}
\opt{A}{The final charge on capacitor $\dfrac{B^2\ell^2CQ_0}{M+CB^2\ell^2}$}
\opt{B}{The maximum (muzzle) speed of the projectile is $\dfrac{2B\ell Q_0}{M+CB^2\ell^2}$}
\opt{C}{The maximum (muzzle) speed of the projectile is $\dfrac{B\ell Q_0}{M+CB^2\ell^2}$}
\opt{D}{The final charge on capacitor $\dfrac{2B^2\ell^2CQ_0}{M+CB^2\ell^2}$}
\end{qoptions}
\qn{4} A small loop of radius $r(\ll R)$, where $R$ = radius of the larger loop that carries a current $i$ is placed on the axis of the bigger loop. If the smaller loop is rotated through $90^\circ$, then the (the resistance of the small loop is $R_0$). [$\omega$ = Angular velocity of loop along the diametrical axis]
\qfig{p1_phy_s1_q04_1}
\begin{qoptions}
\opt{A}{Change in flux is $\dfrac{\pi\mu_0r^2i}{4\sqrt2R}$}
\opt{B}{Average current flowing through smaller loop is $\dfrac{\mu_0r^2\omega i}{2\sqrt2R_0R}$}
\opt{C}{Change in flux is $\dfrac{\pi\mu_0r^2i}{2\sqrt2R}$}
\opt{D}{Average current flowing through smaller loop is $\dfrac{\mu_0r^2\omega i}{2\sqrt2R_0R}$}
\end{qoptions}
\qn{5} \qfig{p1_phy_s1_q05_1}
$R=100\,\Omega$, $L=\dfrac{2}{\pi}$\,H, $C=\dfrac{1}{10\pi}$\,mF\par
In above given AC circuit, $V_R$, $V_L$ and $V_C$ are potential drop across resistance, inductor and capacitor at an instant $t$. Then
\begin{qoptions}
\opt{A}{$V_C=220\sin(100\pi t-\pi/4)$}
\opt{B}{$V_L=440\sin(100\pi t+\pi/4)$}
\opt{C}{$V_C=220\sin(100\pi t-3\pi/4)$}
\opt{D}{$V_L=220\sin(100\pi t-\pi/4)$}
\end{qoptions}
\qn{6} The potential energy between two atoms in a diatomic molecule varies with $x$ as $U=\dfrac{a}{x^{12}}-\dfrac{b}{x^{6}}$, where $a$ and $b$ are positive constants. Consider oscillation of one atom if the other atom is kept stationary.
\qfig{p1_phy_s1_q06_1}
\begin{qoptions}
\opt{A}{Equivalent spring constant $=\left(\dfrac{9b^2}{a}\right)\left(\dfrac{b}{2a}\right)^{1/3}$}
\opt{B}{Equivalent spring constant $=\left(\dfrac{9b^2}{a}\right)\left(\dfrac{2b}{a}\right)^{1/3}$}
\opt{C}{Equilibrium position of particle $=\left(\dfrac{2a}{b}\right)^{1/6}$}
\opt{D}{Equivalent position of particle $=\left(\dfrac{a}{2b}\right)^{1/6}$}
\end{qoptions}
\qn{7} A uniform disc of mass $m$ and radius $R$ is hanging from a rigid support by $N$-number of strings each of length $h$ as shown in the figure. Find the time period of small torsional oscillation about vertical axis passing through the CM of the disc.
\qfig{p1_phy_s1_q07_1}
\begin{qoptions}
\opt{A}{$\pi\left(\dfrac{R}{r}\right)\sqrt{\dfrac{h}{g}}$}
\opt{B}{$2\pi\left(\dfrac{R}{r}\right)\sqrt{\dfrac{2h}{g}}$}
\opt{C}{$2\pi\left(\dfrac{R}{r}\right)\sqrt{\dfrac{h}{g}}$}
\opt{D}{$\pi\left(\dfrac{R}{r}\right)\sqrt{\dfrac{2h}{g}}$}
\end{qoptions}
\qn{8} long thin wooden cylindrical pipe of radius $R$ carrying a uniform surface charge density $\sigma$, is rotating about its axis with an angular velocity $\omega$ that increases slowly with time $t$ as $\omega=kt$, where $k$ is a positive constant. Which of the following statements is/are correct for the region inside the pipe?
\qfig{p1_phy_s1_q08_1}
\begin{qoptions}
\opt{A}{Both the magnetic and the electric fields are uniform and constant.}
\opt{B}{The magnetic field is uniform but not constant and the electric field is constant but not uniform.}
\opt{C}{The magnetic field is constant but not uniform and the electric field is uniform but not constant.}
\opt{D}{Total energy density $u$ due to both the fields varies with time as $u=a+bt^2$, where $a$ and $b$ are positive constants.}
\end{qoptions}
\qn{9} A very long straight conductor and a square conducting frame lie in a plane and are separated from each other as shown in the figure $a=10$ cm. If current in the square frame is increasing a rate of 2 A/s in clockwise direction :-
\qfig{p1_phy_s1_q09_1}
\begin{qoptions}
\opt{A}{Mutual inductance of the two loops is $2\ell n2\times10^{-8}$\,H}
\opt{B}{The current induced in the wire is upwards.}
\opt{C}{The current induced in the wire is downwards.}
\opt{D}{The current induced in the wire is zero.}
\end{qoptions}
\qn{10} An aircraft is flying at a level height in presence the magnetic field of the Earth. If an electric bulb is connected between the two extreme ends of the wings,
\begin{qoptions}
\opt{A}{A voltage will induce across the wings and the bulb will glow.}
\opt{B}{No voltage will induce across the wings but the bulb will glow.}
\opt{C}{A voltage will induce across the wings but the bulb will not glow.}
\opt{D}{No voltage will induce across the wings and the bulb will not glow.}
\end{qoptions}
\sect{Section-II}
\qn{1} A square loop of side $\ell=10$ cm with its sides parallel to $x$ and $y$ axes moves with velocity $\vec v=-7\hat i$ m/s. An outward magnetic field is present which decreases along negative $x$-axis with a gradient $\dfrac{\partial B}{\partial x}=0.1$ T/m and increases with time at a rate of $\dfrac{\partial B}{\partial t}=0.1$ T/s. If the resistance of the loop is $2\times10^{-3}$ ohms, then find the magnitude and direction of induced current in the loop. (in Ampere)
\qfig{p1_phy_s2_q01_1}
\qn{2} Two thin long conducting tapes of which width ``$a$'' are separated through a distance ``$h$'' as shown in the figure. Current $i$ is uniformly distributed over (the area of) each tape. Assuming linear density of the current over its width uniform the inductance of the tape per its unit length is $\dfrac{\mu_0h^{x}}{a^{y}}$. Find value of $x+y$.
\qfig{p1_phy_s2_q02_1}
\qn{3} A resistor $R$ is connected in series with a coil (Resistance not negligible). The system is subjected to an A.C. Supply of peak voltage $V_0$. If the peak voltages across the resistor $R$ and the coils are $V_1$ and $V_2$ respectively, find the power (in watt) dissipated in the coil. (assume $V_1=5$\,V, $V_2=7$\,V, $V_0=10$\,V, $R=13\,\Omega$)
\qn{4} In R--L--C series circuit, we have same current at angular frequencies $\omega_1$ and $\omega_2$. Find the resonant frequency. (Take $\omega_1=4$ units, $\omega_2=9$ units)
\qn{5} A body of mass $m$ hangs from a smooth fixed pulley $P_1$ by the inextensible string fitted with the springs of stiffnesses $k_1$ and $k_3$. The string passes over the smooth light pulley $P_2$ which is connected with another ideal spring of stiffness $k_2$. Find the period of oscillation of the body. (in seconds) assume all springs massless and take $k_1=k_3=\dfrac{k_2}{4}=\pi^2$ SI units, $m=3$ kg
\qfig{p1_phy_s2_q05_1}
\qn{6} A stepped pulley having mass $m$, radius $R$ and radius of gyration about cm is connected with two ideal springs of stiffness $k_1$ and $k_2$ as shown in the figure. If the pulley rolls without sliding, the angular frequency of its oscillation is given by $\sqrt{\dfrac{xk_1r^2+yk_2(R-r)^2}{M(K^2+r^2)}}$. Find value of $x+y$.
\qfig{p1_phy_s2_q06_2}
\qn{7} A simple pendulum has time period $T_0$, say. When it is placed in a liquid the time period decreases. Then we need to move the point of suspension with a vertical acceleration $a$ ($a=2g$) to regain the time period of simple pendulum. The ratio of densities of liquid and the bob is:
\qfig{p1_phy_s2_q07_1}
\qn{8} Period of oscillation of the particle P is $2\pi\sqrt{\dfrac{3m}{xk}}$, then $x$ is equal to
\qfig{p1_phy_s2_q08_1}
\qn{9} The long, horizontal pair of rails shown in the figure is connected using resistance $R$. The distance between the rails is $\ell$, the electrical resistance of the rails is negligible. A conducting wire of mass $m$ and length $\ell$ can slide without friction on the pair of rails, in a vertical, homogeneous magnetic field of induction $B$. A force of magnitude $F_0$ is exerted for sufficiently long time onto the conducting wire, so that the speed of the wire becomes nearly constant. The force $F_0$ is now removed at a certain point P. The distance the conducting wire cover on rails from point P before stopping is $(316+n)$ meters. Find `$n$'. (Given : $F_0=20$\,N, $m=1.6$ gm, $R=0.01\,\Omega$, $\ell=10$ cm, $B=0.1$ T)
\qfig{p1_phy_s2_q09_1}
\qn{10} A weight rigid rod with a small iron bob at the end is hinged at point A to the wall so that it can rotate in all directions. The rod is kept in the horizontal position by a vertical inextensible string of length 20\,cm, fixed at its midpoint. The bob is displaced slightly perpendicular to the plane of the rod and string. Find period of small oscillations of the system. If time period is $\dfrac{2\pi}{x}$ write $x$.
\qfig{p1_phy_s2_q10_1}
\partsec{2}{CHEMISTRY}
\sect{Section-I}
\qn{1} An oxime of compound\qfig{p1_chem_s1_q01_1}whenever treated with H$_2$SO$_4/\Delta$ gives a compound P, which on hydrolysis produces Q and R. If Q is an aliphatic primary amine, then choose the correct option(s)?
\begin{qoptions}
\opt{A}{R is soluble in hot aqueous NaOH}
\opt{B}{R gives brisk effervence of CO$_2$ with aqueous NaHCO$_3$}
\opt{C}{Q can give foul smelling compound \par\smallskip\qfigi{p1_chem_s1_q01_2}}
\opt{D}{Q on Hoffmann exhaustive methylation gives product which has very high solubility in water}
\end{qoptions}
\qn{2} CH$_3-$CH$_2-$CH$_2-$NH$_2$ $\xrightarrow{\text{N}_2\text{O}_3}$ The product(s) formed is/are
\begin{qoptions}
\opt{A}{\qfigi{p1_chem_s1_q02_A}}
\opt{B}{CH$_3-$CH$_2-$CH$_2-$OH}
\opt{C}{CH$_3-$CH$=$CH$_2$}
\opt{D}{CH$_3-$CH$_2-$C$\equiv$N}
\end{qoptions}
\qn{3} By reacting with which of the following reagents CH$_3$NH$_2$ will form a solid derivative?
\begin{qoptions}
\opt{A}{PhSO$_2$Cl}
\opt{B}{(COOEt)$_2$}
\opt{C}{H$_2$SO$_4$}
\opt{D}{CHCl$_3$ + KOH}
\end{qoptions}
\qn{4} A plot of $[\mathrm A]_t$ vs $t$ (where, $[\mathrm A]_t$ is the concentration of the reactant at time $t$) for a reaction is linear with a negative slope of 0.0231 sec$^{-1}$. Then
\begin{qoptions}
\opt{A}{The initial concentration of the reactant will reduce to half in 30 seconds.}
\opt{B}{A plot of $t_{1/2}$ vs concentration of reactant will give a straight line passing through origin.}
\opt{C}{The reaction will be 90\% complete in 100 seconds.}
\opt{D}{The line $\ln[\mathrm A]_t$ vs $t$ plot will be steeper at higher temperature.}
\end{qoptions}
\qn{5} For a first order reaction X(g) $\to$ 2Y(g) + Z(g), initially only 3 moles of X(g) was taken in 10\,L rigid closed container at 300\,K. After 10 seconds the reaction was stopped forcefully and a pinhole was made in the container at the same temperature. The molar ratio of X(g) and Y(g) gases coming out is 1:2. The \textbf{CORRECT} option(s) is/are:\par
(Given molar masses of X, Y and Z are 16, 4 and 8\,g/mole respectively $\ell$n2 = 0.7 and $\ell$n3 = 1.1, $R=0.08$ atm L$^{-1}$K$^{-1}$mol$^{-1}$)
\begin{qoptions}
\opt{A}{The total pressure inside the container when reaction was just stopped is 12 atm.}
\opt{B}{The rate constant of reaction is $4\times10^{-2}$ sec$^{-1}$}
\opt{C}{Moles of Z just after reaction was stopped forcefully is 1.0.}
\opt{D}{The rate of formation of Y after 10 second is 0.4 mol L$^{-1}$ sec$^{-1}$.}
\end{qoptions}
\qn{6} ACl$_3$ is an unstable compound and gradually changing into A$_2$O$_5$ by its reaction with atmospheric oxygen, according to the following partial reaction:\par
2ACl$_3\to$ A$_2$O$_5$\par
The above reaction follows second order kinetics. The reaction is started by taking 1 mole of ACl$_3$ in a container of volume 1\,L and allowed to stand for some time. After 10\,min, the reaction mixture requires 250 mL of 0.8\,M acidified KMnO$_4$ solution; where ACl$_3$ is oxidized into AO$_3^-$ and assume A$_2$O$_5$ do not react with KMnO$_4$. Now, choose the correct statement(s) of the following :
\begin{qoptions}
\opt{A}{Half-life time of the reaction is 10 minutes}
\opt{B}{Rate constant of the reaction is 0.05 $\ell$\,mol$^{-1}$min.$^{-1}$}
\opt{C}{Rate constant of the reaction is 0.10 $\ell$\,mol$^{-1}$min.$^{-1}$}
\opt{D}{After 20 min from the start, 22.33\% of ACl$_3$ is left unreached}
\end{qoptions}
\qn{7} The observed rate of a chemical reaction is substantially lower than the collision frequency. One or more of the following statements is/are true to account for this fact.
\begin{qoptions}
\opt{A}{The reactants do not have the required energy}
\opt{B}{The partners do not collide in the proper orientation.}
\opt{C}{Collision complex exists for a very short time.}
\opt{D}{Collision frequency overestimates the number of effective collisions.}
\end{qoptions}
\qn{8} The correct statement(s) regarding the halides of boron is/are
\begin{qoptions}
\opt{A}{BI$_3$ undergo the easiest hydrolysis than the other halides}
\opt{B}{BF$_3$ + BCl$_3\to$ BF$_2$Cl + BFCl$_2$}
\opt{C}{Inorganic benzene is formed when BCl$_3$ reacts with NH$_4$Cl followed by reaction with NaBH$_4$}
\opt{D}{Lewis acidity order: BBr$_3>$ BCl$_3>$ BF$_3$}
\end{qoptions}
\qfullbegin
\qn{9} Identify A, B, C and D in the following reaction sequence.
\qfig{p1_chem_s1_q09_1}
\qfullsplit
\begin{qoptions}
\opt{A}{The compound A is (CH$_3$)$_2$SiCl$_2$}
\opt{B}{The compound C and D are SiO$_2$ and SiC respectively}
\opt{C}{The compound B is (CH$_3$)$_2$Si(OH)$_2$}
\opt{D}{The compound B is (CH$_3$)$_2$Si$=$O}
\end{qoptions}
\qfullend
\qn{10} Inorganic benzene surpasses organic benzene in which property:
\begin{qoptions}
\opt{A}{Melting point}
\opt{B}{Edge length of its hexagon}
\opt{C}{Molar mass}
\opt{D}{Total number of electrons in the molecule}
\end{qoptions}
\sect{Section-II}
\qn{1} A solid compound `G', of formula C$_{15}$H$_{15}$ON, was insoluble in water, dilute HCl, of dilute NaOH. After prolonged heating of `G' with aqueous NaOH, a liquid, `H', was observed floating on the surface of the alkaline mixture. `H' did not solidify upon cooling to room temperature; it was steam-distilled and separated. Acidification of the alkaline mixture with hydrochloric acid caused precipitation of a white solid, `I'. `I' is a paramethyl benzene derivative.\par
Compound `H' was soluble in dilute HCl, and reacted with benzenesulfonyl chloride and excess KOH to give a base-insoluble solid, `J'.\par
Compound `I', m.p. 180$^\circ$C, was soluble in aqueous NaHCO$_3$, and contained no nitrogen.\par
The number of Carbon atoms in `I' will be `\textbf{a}'\par
The $\dfrac{\text{Molecular weight of `H'}}{53.5}$ will be `\textbf{b}'\par
Then the value of $(\mathbf{a}+\mathbf{b})$ is
\qn{2} An aromatic compound P(C$_8$H$_{10}$) is heated with alkaline KMnO$_4$ to give a product which on heating with aqueous ammonia solution gives `Q'. Compound `Q' on heating with alkaline sodium hypochlorite gives a salt which on acidification gives product R (C$_7$H$_7$O$_2$N). Compound `R' is treated with NaNO$_2$ \& HCl and then with a base to produce a dipolar ion `S' which on heating produces an intermediate (T) which dimerises to produce the final product `U'. What is the degree of unsaturation of the compound (U)?
\qfullbegin
\qn{3} \qfig{p1_chem_s2_q03_1}
\qfullsplit
\qfig{p1_chem_s2_q03_2}
In the above reaction sequence molar mass of U is \ldots
\qfullend
\qn{4} The number of following amines, which are relatively more basic than benzyl amine is/are $\mathbf a$ and the amines which can be easily prepared by Gabriel's phthalimide synthesis is/are $\mathbf b$ then the value of $\mathbf{a/b}$ would be :
\qfig{p1_chem_s2_q04_1}
\qn{5} If certain decomposition reaction found to obey equation in term of partial pressure $(\mathbf P)$ of reactant as :\par $\dfrac{a-P^2}{P^2}=bkt$, what is order of reaction if $\mathbf a$, $\mathbf b$ and $\mathbf k$ are constant ($t$ represents time) ?
\qn{6} An optically active substance, A, decomposes into optically active substance `B' and `C' as :\par
A $\xrightarrow{K=0.001\ \mathrm{min}^{-1}}$ 2B + C\par
The specific rotations of A, B and C are $+40^\circ$, $10^\circ$ and $-30^\circ$ per mole, respectively. If initially A and C were present in 4 : 3 mole ratio, the time (in min), after which the sample becomes optically inactive, is (ln 2 = 0.7, ln 5 = 1.6, ln 7 = 2.0, ln 13 = 2.5)
\qn{7} A 1.0\,mL solution containing 0.240\,$\mu$Ci of tritium is injected into a Dog's blood stream. After a period of time to allow the isotope to be dispersed, a 1.00\,mL sample of blood is drawn. The radioactivity of this sample is found to be $3.2\times10^{4}$\,$\mu$Ci/mL. The total volume of blood in the Dog in litre is \ldots?
\qn{8} A graph between $\log t_{1/2}$ and $\log a$ (abscissa), $a$ being the initial concentration of $\mathbf A$ in the reaction. A $\to$ Product is depicted in figure.
\qfig{p1_chem_s2_q08_1}
The rate law is $-\dfrac{d[\mathrm A]}{dt}=k[\mathrm A]^{\mathbf n}$\par
The value of $\mathbf n$ will be
\qn{9} Total number of OH groups present on all boron atoms in colemanite Ca$_2$B$_6$O$_{11}\cdot$5H$_2$O is
\qn{10} When one mole of K$_4\left[\mathrm{Fe(CN)}_6\right]$ is heated with conc. H$_2$SO$_4$, the evolved gas during reaction is reacted with excess of I$_2$O$_5$. Find how many moles of I$_2$ will be liberated during the reaction
\partsec{3}{MATHEMATICS}
\sect{Section-I}
\qn{1} $\displaystyle\int_0^{\pi/2}\tan^3x\cos^5x\cos7x\,dx$
\begin{qoptions}
\opt{A}{$\dfrac18\displaystyle\int_{\pi}^{0}\sin^2x\sin3x\,dx$}
\opt{B}{$\dfrac18\displaystyle\int_{\pi/2}^{0}\sin^2x\sin3x\,dx$}
\opt{C}{$\dfrac1{30}$}
\opt{D}{$\dfrac1{60}$}
\end{qoptions}
\qn{2} Let $I_n=\displaystyle\int_0^{n}2^x\left(1+\{g(x)\}+\left\{\sqrt{g(x)}\right\}\right)dx$, $n\in N$, where $g(x)=\sin^2(\pi x)$ ($\{x\}$ represents fractional part of $x$), then
\begin{qoptions}
\opt{A}{$I_1>\log_2e$}
\opt{B}{$I_1<\log_2e$}
\opt{C}{$I_4=15I_1$}
\opt{D}{$I_4=16I_1$}
\end{qoptions}
\qn{3} If $f(\alpha)=\displaystyle\lim_{n\to\infty}\frac1n\sum_{r=0}^{n\alpha}\frac{2r}{\sqrt{(3n\alpha-2r)(n\alpha+2r)}}$, then
\begin{qoptions}
\opt{A}{$f(2)=\dfrac{\pi}{3}$}
\opt{B}{$f(\alpha)=\dfrac{\alpha}{2}\displaystyle\int_0^{\alpha}\dfrac{dx}{\sqrt{\alpha^2-\left(x-\frac{\alpha}{2}\right)^2}}$}
\opt{C}{$f(\alpha)$ is a constant function}
\opt{D}{$f'(\alpha)=\dfrac{\pi}{6}$}
\end{qoptions}
\qn{4} If $\displaystyle\int_0^{x}\frac{(mt\cos4t-\ell\sin4t)\,dt}{t^2}=\frac{\ell\sin4x}{x}-1$, where $\left(0<x<\dfrac{\pi}{4}\right)$, then $(\ell+m)$ equals
\begin{qoptions}
\opt{A}{$\dfrac54$}
\opt{B}{$\dfrac74$}
\opt{C}{$\dfrac23$}
\opt{D}{$\dfrac15$}
\end{qoptions}
\qn{5} Let $f(x)$ is a real valued function defined by $f(x)=x^2+x^2\displaystyle\int_{-1}^{1}tf(t)\,dt+x^3\int_{-1}^{1}f(t)\,dt$, then which of the following hold(s) good?
\begin{qoptions}
\opt{A}{$\displaystyle\int_{-1}^{1}tf(t)=\dfrac{10}{11}$}
\opt{B}{$f(1)+f(-1)=\dfrac{30}{11}$}
\opt{C}{$\displaystyle\int_{-1}^{1}tf(t)\,dt>\int_{-1}^{1}f(t)\,dt$}
\opt{D}{$f(1)-f(-1)=\dfrac{20}{11}$}
\end{qoptions}
\qn{6} Which of the following is true
\begin{qoptions}
\opt{A}{$\displaystyle\int_{-\pi/2}^{\pi/2}\dfrac{x^2\cos x}{1+e^x}\,dx=\dfrac{\pi^2}{4}-2$}
\opt{B}{Area bounded by two branches of curve $(y-x)^2=x^3$ and $x=1$ is equal to $\dfrac45$}
\opt{C}{If $f(x)=\dfrac{x^3+x}{2}$, $f:R\to R$ \& $g(x)=f^{-1}(x)$ then area bounded between $f(x)$ and $g(x)$ is equal to $\dfrac14$}
\opt{D}{Area bounded by $y^2=4x$ and normal to $y^2=4x$ with slope $-1$ is equal to $\dfrac{64}{3}$}
\end{qoptions}
\qn{7} Which of the following is/are true
\begin{qoptions}
\opt{A}{$\displaystyle\int_0^{1/2}\cos x\,dx>\dfrac{23}{48}$}
\opt{B}{$\displaystyle\int_0^{1/2}x\sin x\,dx<\dfrac1{24}$}
\opt{C}{$\displaystyle\int_0^{1/2}x\cos x\,dx>\dfrac{15}{128}$}
\opt{D}{$\displaystyle\int_0^{1/2}\sin x\,dx<\dfrac{47}{384}$}
\end{qoptions}
\qn{8} Let $f(x)=\begin{cases}[x], & 1\le x<3\\ 2x, & 3\le x<4\\ 5, & 4\le x<6\end{cases}$ where $[\,.\,]$ denotes greatest integer function and $g(x)=\displaystyle\int_1^{x}f(t)\,dt$ then
\begin{qoptions}
\opt{A}{$\displaystyle\int_1^3g(t)\,dt=\dfrac52$}
\opt{B}{$\displaystyle\int_1^3g(t)\,dt=3$}
\opt{C}{$\displaystyle\int_3^5g(t)\,dt=\dfrac{113}{6}$}
\opt{D}{$g(5)=15$}
\end{qoptions}
\qn{9} If $f(x)=\max\{|x|,|x-1|,|x-2|,\ldots,|x-10|\}$ and $g(x)=\min\{|x|,|x-1|,|x-2|,\ldots|x-10|\}$ then
\begin{qoptions}
\opt{A}{Area bounded by $y=f(x)$, $x$-axis, $x=0$ and $x=10$ is 75 unit}
\opt{B}{Area bounded by $y=g(x)$, $y=1$, $x=0$ and $x=10$ is 7.5}
\opt{C}{Area bounded by $y=f(x)$, $y=g(x)$, $x=0$ and $x=10$ is 72.5}
\opt{D}{Area bounded by $y=f(x)$, $y=g(x)$, $x=0$ and $x=10$ is 70}
\end{qoptions}
\qn{10} Let $f$ be a differentiable function on $R$ and satisfy the $\displaystyle\int_0^{x}f(t)\,dt+\int_0^{x}tf(x-t)\,dt=-1+e^{-x}\ \forall x\in R$. Which of the following is/are correct
\begin{qoptions}
\opt{A}{$f(0)+f'(0)+f(2)=1+e^{-2}$}
\opt{B}{$f(0)+f'(0)+f(2)=2+e^{-2}$}
\opt{C}{Area enclosed by $y=f(x)$, $y=0$, $x=0$ and $x=2$ is $\dfrac1e$}
\opt{D}{Area enclosed by $y=f(x)$, $y=0$, $x=0$ and $x=2$ is $\dfrac2e-\dfrac2{e^2}$}
\end{qoptions}
\sect{Section-II}
\qn{1} Let $\sin\beta\cos x-\tan\alpha\sin x+2\sin2x=0$ and
{\footnotesize\[f(x)=\int_0^{x}\frac{\begin{array}{c}(\cos3x-\cos x)(\tan\alpha-1)^2\\+(\sin3x+\sin x)(1-\sin\beta)^2\\-\sin2x\left(\frac{\tan\alpha-\sin\beta}{\sqrt2}\right)^2\end{array}}{\cos x\sin^2\beta-\sin x\tan^2\alpha+2\sin2x}\,dx,\]}
where $\cos x\sin^2\beta-\sin x\tan^2\alpha+2\sin2x\neq0$. If $M$ is maximum value of $(f(x)+\cos2x)^2$, $x\in\left(0,\dfrac{\pi}{2}\right)$, then value of $M$ is
\qn{2} If $\displaystyle\int_1^{2}(x^3+3)(x^3-1)\,dx=\lambda\int_2^{9}x\left(2-(x-1)^{1/3}\right)dx$ then find $\lambda$.
\qn{3} Let $f(x)=\displaystyle\lim_{n\to\infty}\sum_{r=1}^{n}\frac{x}{n}e^{\frac{rx}{n}}\sin\left(\frac{n-r}{n}x\right)$. If $f''(x)-f(x)=a\sin x+b\cos x$, then find the value of $\left(\dfrac{a+b}{5}\right)$.
\qn{4} If $\displaystyle\lim_{x\to0}\dfrac{\displaystyle\int_0^{x}\dfrac{e^{x-t}\,dt}{\sqrt{t^2+x^2+1-2tx}}}{x(e^x-1)}=\lambda$, then find the value of $\lambda$
\qn{5} Let $u=\displaystyle\int_0^{\pi/2}\cos\left(\dfrac{2\pi}{3}\sin^2x\right)dx$ and $v=\displaystyle\int_0^{\pi/2}\cos\left(\dfrac{\pi}{3}\sin x\right)dx$, then $\dfrac uv$ equals
\qn{6} Let $f(x)$, $(x\ge1)$ be a differentiable function satisfying $f(x)=\ell n^2x-\displaystyle\int_1^{e}\dfrac{f(t)}{t}\,dt$, then area bounded by tangent line of $y=f(x)$ at $x=e$, the curve $y=f(x)$ and the line $x=1$ is $A$, then $[A]$ is equal to (where $[\,.\,]$ denotes greatest integer function)
\qn{7} A function $y=f(x)$ satisfies $xf'(x)-2f(x)=x^4f^2(x)$, $\forall\ x>0$ and $f(1)=-6$. Find value of $f'\left(3^{1/5}\right)$
\qn{8} $f(x)=\left(e^{-x}+e^{x}\right)\cos x-2x-\displaystyle\int_0^{x}(x-t)f'(t)\,dt$. If $y=f(x)$ satisfies above equation and $f\left(\dfrac{\pi}{2}\right)$ is equal to $Ae^{\pi/2}+Be^{-\pi/2}$, then $|5(A+B)|$ is equal to (where $A,B\in Q$)
\qn{9} Let $f(x)=\begin{cases}2\ell n(-x-1), & -e-1\le x\le-2\\[2pt] \frac12(4-x^2), & -2<x<0\\[2pt] 0, & x=0\\[2pt] \frac12(x^2-4), & 0<x<2\\[2pt] 2\ell n(x-1), & 2\le x\le e+1\end{cases}$\par
Also, {\small$g(x)=\begin{cases}\min\{f(t):-e-1\le t\le x\}, & -e-1\le x<0\\[2pt] \max\{f(t):0\le t\le x\}, & 0\le x\le e+1\end{cases}$}\par
then value of $\displaystyle\int_{-e-1}^{e+1}g(x)\,dx$ is equal to $A$, then value of $[A]$ is, (where $[\,.\,]$ denotes greatest integer function)
\qn{10} $f(x)$ and $g(x)$ are two differentiable function defined on $R$, satisfying $f(x)=\dfrac{x^3}{2}+1-x\displaystyle\int_0^{x}g(t)\,dt$ and $g(x)=x-\displaystyle\int_0^{1}f(t)\,dt$, then the area bounded by $y=f(x)$, $y=g(x)$, $x=0$ and $x=2$ is equal to
\end{multicols}
\end{document}
